\chapter{Categories}

\begin{center}\itshape
What do a free group, the abelianisation of a group, the double dual of a vector space and currying a function of two variables have in common? Each is the best possible
solution of a problem about maps. Category theory is the language in which such statements become precise, and its first theorem, Yoneda's lemma, says that an object is
known completely by the maps into it.
\end{center}

In the early 1940s Saunders Mac Lane gave an invited lecture at the University of Michigan on group extensions: given two abelian groups $A$ and $B$, classify the groups $E$
containing $A$ with $E/A\cong B$. His answer was a group he called $\operatorname{Ext}(B,A)$. In the audience sat Samuel Eilenberg, a young topologist who had left Poland in 1939,
and he saw at once that Mac Lane's group answered a question of Norman Steenrod about the homology of the complement of a strange space, the $p$-adic solenoid. A long
collaboration began. To explain why their isomorphisms were not accidents of a choice of basis, they needed a word for what
mathematicians had always called ``natural'' without defining it. In 1945 they published ``General theory of natural equivalences''. Mac Lane later summed up the logic:
categories were defined in order to define functors, and functors in order to define natural transformations.

For a decade the subject was regarded as a convenient bookkeeping, ``general abstract nonsense'' as its own users fondly called it. Then came a conversation in Paris. Nobuo
Yoneda, a young Japanese mathematician travelling in Europe, met Mac Lane in a café of the Gare du Nord; the interview went on in Yoneda's train until it left the station,
and what Yoneda explained there is the lemma that bears his name. In 1958 Daniel Kan defined adjoint functors, and the bookkeeping turned out to contain theorems.

This chapter sets up the language that the rest of this tome speaks: categories, functors and natural transformations (Tome~2, Chapter~12 promised them), universal
properties, the Yoneda lemma, and adjunctions. Every definition comes with examples from the earlier tomes, because the point of the language is that they were already
there.

\section{Categories}

\begin{definition}\label{def:category}
A \emph{category} $\mathcal C$ consists of a class of \emph{objects}; for any two objects $A,B$ a set $\operatorname{Hom}_{\mathcal C}(A,B)$ of \emph{morphisms} (or \emph{arrows}) from $A$ to $B$,
written $f\colon A\to B$; and for any three objects a \emph{composition} $\operatorname{Hom}(B,C)\times\operatorname{Hom}(A,B)\to\operatorname{Hom}(A,C)$, $(g,f)\mapsto g\circ f$, such that
\begin{enumerate}[label=(\roman*)]
\item composition is associative: $h\circ(g\circ f)=(h\circ g)\circ f$;
\item every object $A$ has an \emph{identity} $\mathrm{id}_A\in\operatorname{Hom}(A,A)$ with $f\circ\mathrm{id}_A=f$ and $\mathrm{id}_B\circ f=f$ for all $f\colon A\to B$.
\end{enumerate}
\end{definition}

The objects of $\mathbf{Set}$ do not form a set (there is no set of all sets, Tome~6, Chapter~1), which is why the definition says ``class''. We shall never need more than this
word: every argument of this chapter takes place inside finitely many $\operatorname{Hom}$-sets.

\begin{example}\label{ex:categories}
\leavevmode
\begin{enumerate}[label=(\alph*)]
\item The \emph{concrete} categories: $\mathbf{Set}$ (sets and maps), $\mathbf{Grp}$ (groups and homomorphisms), $\mathbf{Ab}$ (abelian groups), $\mathbf{Ring}$ (rings with $1$ and
homomorphisms preserving $1$), $\mathbf{CRing}$ (commutative rings), $R\text{-}\mathbf{Mod}$ (modules over a ring $R$ and $R$-linear maps), $\mathbf{Vect}_K$ (vector spaces over a field $K$),
$\mathbf{Top}$ (topological spaces and continuous maps), $\mathbf{Top}_*$ (spaces with a base point and maps preserving it). Composition is composition of maps.
\item A \emph{poset} $(P,\le)$ is a category with one morphism $x\to y$ when $x\le y$ and none otherwise; reflexivity gives the identities and transitivity the composition.
\item A \emph{monoid} $M$ is a category with one object $\bullet$ and $\operatorname{Hom}(\bullet,\bullet)=M$, composition being the product of $M$. A group is a one-object category in which every
morphism has an inverse.
\item $\mathbf{Mat}_K$ has the natural numbers as objects, the $m\times n$ matrices as morphisms $n\to m$, and the product of matrices as composition. Nothing in it is a set with
structure: a morphism is not a map between anything.
\item The \emph{opposite} category $\mathcal C^{\mathrm{op}}$ has the objects of $\mathcal C$, $\operatorname{Hom}_{\mathcal C^{\mathrm{op}}}(A,B)=\operatorname{Hom}_{\mathcal C}(B,A)$, and composition reversed.
\end{enumerate}
\end{example}

A morphism $f\colon A\to B$ is an \emph{isomorphism} if there is $g\colon B\to A$ with $g\circ f=\mathrm{id}_A$ and $f\circ g=\mathrm{id}_B$. Such a $g$ is unique (if $g'$ is another,
$g'=g'\circ f\circ g=g$), and is written $f^{-1}$. In $\mathbf{Set}$ the isomorphisms are the bijections, in $\mathbf{Grp}$ the group isomorphisms, in $\mathbf{Top}$ the homeomorphisms,
which are not the same as the continuous bijections. A morphism $f$ is a \emph{monomorphism} if $f\circ g=f\circ h$ implies $g=h$, and an \emph{epimorphism} if $g\circ f=h\circ f$ implies $g=h$.
In $\mathbf{Set}$ these are the injections and the surjections, but Exercise~\ref{ex:zq} shows that a morphism of rings can be both without being an isomorphism.

\begin{think}
In the poset category, when is a morphism an isomorphism? What is an isomorphism in $\mathbf{Mat}_K$, and why are there no isomorphisms $m\to n$ for $m\ne n$?
\end{think}

\section{Functors}

\begin{definition}\label{def:functor}
A \emph{functor} $F\colon\mathcal C\to\mathcal D$ assigns to each object $A$ of $\mathcal C$ an object $F(A)$ of $\mathcal D$, and to each morphism $f\colon A\to B$ a morphism $F(f)\colon F(A)\to F(B)$, so that
\[F(g\circ f)=F(g)\circ F(f)\qquad\text{and}\qquad F(\mathrm{id}_A)=\mathrm{id}_{F(A)}.\]
A \emph{contravariant functor} from $\mathcal C$ to $\mathcal D$ is a functor $\mathcal C^{\mathrm{op}}\to\mathcal D$: it reverses arrows, $F(g\circ f)=F(f)\circ F(g)$.
\end{definition}

\begin{example}\label{ex:functors}
\leavevmode
\begin{enumerate}[label=(\alph*)]
\item \emph{Forgetful functors} $U\colon\mathbf{Grp}\to\mathbf{Set}$, $\mathbf{Ring}\to\mathbf{Ab}$, $\mathbf{Top}\to\mathbf{Set}$, which forget part of the structure.
\item \emph{Abelianisation} $\mathbf{Grp}\to\mathbf{Ab}$, $G\mapsto G^{\mathrm{ab}}=G/G'$, where $G'$ is the subgroup generated by the commutators $aba^{-1}b^{-1}$ (Tome~2, Exercise~2.4). A homomorphism $f\colon G\to H$ maps commutators to commutators, hence $G'$ into $H'$, and induces $f^{\mathrm{ab}}\colon gG'\mapsto f(g)H'$.
\item The units $R\mapsto R^\times$ and $R\mapsto\mathrm{GL}_n(R)$ are functors $\mathbf{CRing}\to\mathbf{Grp}$: a ring homomorphism, applied to each entry, maps invertible matrices to invertible
matrices.
\item The fundamental group $\pi_1\colon\mathbf{Top}_*\to\mathbf{Grp}$ (Tome~8, Chapter~5), and the homology groups $H_n\colon\mathbf{Top}\to\mathbf{Ab}$ (Tome~8, Chapter~6). That $H_n(g\circ f)=H_n(g)\circ H_n(f)$
was proved there; it is exactly the statement that $H_n$ is a functor.
\item Duality $V\mapsto V^*=\operatorname{Hom}(V,K)$, $f\mapsto f^*=(\varphi\mapsto\varphi\circ f)$, is a contravariant functor $\mathbf{Vect}_K\to\mathbf{Vect}_K$: $(g\circ f)^*=f^*\circ g^*$.
\item For an object $A$ of any category, the functor $h^A=\operatorname{Hom}(A,-)$ from $\mathcal C$ to $\mathbf{Set}$ sends $B$ to $\operatorname{Hom}(A,B)$ and $g\colon B\to C$ to $g\circ-$. Likewise
the contravariant functor $h_A=\operatorname{Hom}(-,A)$ sends $g\colon B\to C$ to the map $-\circ g$ from $\operatorname{Hom}(C,A)$ to $\operatorname{Hom}(B,A)$.
\item A functor between two posets is an order-preserving map; a functor between two groups, seen as one-object categories, is a homomorphism.
\end{enumerate}
\end{example}

\begin{proposition}\label{prop:functoriso}
A functor maps isomorphisms to isomorphisms: $F(f^{-1})=F(f)^{-1}$.
\end{proposition}

\begin{proof}
$F(f^{-1})\circ F(f)=F(f^{-1}\circ f)=F(\mathrm{id}_A)=\mathrm{id}_{F(A)}$, and likewise in the other order.
\end{proof}

This one-line statement is the engine of algebraic topology: to prove that two spaces are not homeomorphic, find a functor that takes them to non-isomorphic groups. The same idea
proves things about maps.

\begin{example}[Brouwer, functorially]\label{ex:brouwer}
There is no continuous retraction $r\colon D^2\to S^1$ of the closed disc onto its boundary circle. If there were, with $i\colon S^1\to D^2$ the inclusion, $r\circ i=\mathrm{id}_{S^1}$, and applying
the functor $H_1$ gives $H_1(r)\circ H_1(i)=\mathrm{id}_{H_1(S^1)}$. But $H_1(i)$ goes through $H_1(D^2)=0$ while $H_1(S^1)\cong\Z$ (Tome~8, Chapter~6): the identity of $\Z$ would factor through
the zero group. As Tome~8 showed, this gives Brouwer's fixed point theorem, since a map $D^2\to D^2$ without fixed point yields a retraction by pushing each point away from its image to the boundary.
\end{example}

\section{Natural transformations}

\begin{definition}\label{def:nat}
Let $F,G\colon\mathcal C\to\mathcal D$ be functors. A \emph{natural transformation} $\eta\colon F\Rightarrow G$ is a family of morphisms $\eta_A\colon F(A)\to G(A)$, one for each object $A$, such that
for every $f\colon A\to B$ the square of Figure~\ref{fig:universal} commutes:
\[G(f)\circ\eta_A=\eta_B\circ F(f).\]
It is a \emph{natural isomorphism} if every $\eta_A$ is an isomorphism; then $(\eta_A^{-1})$ is a natural transformation $G\Rightarrow F$, as one sees by composing the square with $\eta_A^{-1}$ and $\eta_B^{-1}$.
\end{definition}

A natural transformation is a construction that works the same way for all objects at once, without looking inside them. Here are two that the reader has met.

\begin{example}\label{ex:nattrans}
\leavevmode
\begin{enumerate}[label=(\alph*)]
\item The quotient maps $\pi_G\colon G\to G^{\mathrm{ab}}$ form a natural transformation from the identity functor of $\mathbf{Grp}$ to abelianisation (followed by the inclusion of $\mathbf{Ab}$ in
$\mathbf{Grp}$): $f^{\mathrm{ab}}(\pi_G(g))=f(g)H'=\pi_H(f(g))$.
\item The determinant $\det_R\colon\mathrm{GL}_n(R)\to R^\times$ is natural in the commutative ring $R$: for a ring homomorphism $\varphi$, $\det(\varphi(M))=\varphi(\det M)$, because the
determinant is a polynomial with integer coefficients in the entries.
\end{enumerate}
\end{example}

\begin{think}
Let $f\colon H\to K$ be a morphism. Check that the maps $-\circ f\colon\operatorname{Hom}(K,X)\to\operatorname{Hom}(H,X)$, one for each object $X$, form a natural transformation $h^K\Rightarrow h^H$ between functors
$\mathcal C\to\mathbf{Set}$. In which direction does it go, compared with $f$?
\end{think}

The first precise theorem of the subject is the one Eilenberg and Mac Lane wrote it for. Tome~2, Theorem~A.8 showed that the map $\iota_V\colon V\to V^{**}$, $\iota_V(v)(\varphi)=\varphi(v)$,
is an isomorphism for finite-dimensional $V$, and remarked that $V\cong V^*$ holds too but only through a choice of basis. The language of this chapter turns this remark into a theorem.

\begin{theorem}[the double dual is natural]\label{thm:bidualnat}
On the category $\mathbf{FinVect}_K$ of finite-dimensional vector spaces over $K$, the maps $\iota_V$ form a natural isomorphism from the identity functor to the functor $V\mapsto V^{**}$,
$f\mapsto f^{**}$.
\end{theorem}

\begin{proof}
Each $\iota_V$ is an isomorphism by Tome~2, Theorem~A.8. For $f\colon V\to W$ linear, $v\in V$ and $\psi\in W^*$,
\[f^{**}(\iota_V(v))(\psi)=\iota_V(v)(f^*\psi)=\iota_V(v)(\psi\circ f)=\psi(f(v))=\iota_W(f(v))(\psi),\]
so $f^{**}\circ\iota_V=\iota_W\circ f$.
\end{proof}

For the dual itself, the question needs care, because $V\mapsto V^*$ is contravariant and the identity functor covariant. Both become covariant on the category $\mathbf{Iso}_K$ whose objects are
finite-dimensional spaces and whose morphisms are the linear isomorphisms only, if the dual acts on $f$ by $(f^*)^{-1}$. A natural isomorphism $V\cong V^*$ would then be a family of
isomorphisms $\eta_V\colon V\to V^*$ with $(f^*)^{-1}\circ\eta_V=\eta_W\circ f$, that is,
\begin{equation}\label{eq:dualnat}
\eta_V=f^*\circ\eta_W\circ f
\end{equation}
for every isomorphism $f\colon V\to W$.

\begin{theorem}[$V$ is not naturally its dual]\label{thm:notnatural}
If the field $K$ has more than three elements, there is no family of isomorphisms $\eta_V\colon V\to V^*$, for $V$ in $\mathbf{Iso}_K$, satisfying \eqref{eq:dualnat}. Indeed \eqref{eq:dualnat} already fails
for every non-zero $V$ and the maps $f=\lambda\,\mathrm{id}_V$.
\end{theorem}

\begin{proof}
Let $V\ne0$ and $\lambda\in K^\times$. For $f=\lambda\,\mathrm{id}_V$ we have $f^*=\lambda\,\mathrm{id}_{V^*}$, so \eqref{eq:dualnat} with $W=V$ reads $\eta_V=\lambda^2\eta_V$. Since $\eta_V$ is an isomorphism
of a non-zero space, $\eta_V\ne0$, hence $\lambda^2=1$ for every $\lambda\in K^\times$. A polynomial of degree~$2$ has at most two roots in a field, so $|K^\times|\le2$ and $|K|\le3$, a contradiction.
\end{proof}

The restriction on $K$ is not a weakness of the proof: over $\mathbb F_2$ the planes do have a natural isomorphism with their duals (Exercise~\ref{ex:f2plane}). Theorem~\ref{thm:bidualnat} and
Theorem~\ref{thm:notnatural} together are the precise form of an intuition every student has: a vector is a measurement of measurements, while turning a vector into a measurement needs extra
structure, such as an inner product.

\section{Universal properties}

Many objects of algebra are defined by a construction, then characterised by what maps out of them or into them do. The second description is the one that matters.

\begin{definition}\label{def:initial}
An object $I$ of $\mathcal C$ is \emph{initial} if for every object $X$ there is exactly one morphism $I\to X$, and \emph{terminal} if for every $X$ there is exactly one morphism $X\to I$.
\end{definition}

In $\mathbf{Set}$ the empty set is initial and every one-point set terminal; in $\mathbf{Grp}$ the trivial group is both; in $\mathbf{Ring}$, $\Z$ is initial (the only ring homomorphism $\Z\to R$
sends $n$ to $n\cdot1$) and the zero ring terminal; in a poset an initial object is a least element.

\begin{proposition}\label{prop:initialunique}
Two initial objects $I,I'$ are isomorphic, by a unique isomorphism. The same holds for terminal objects.
\end{proposition}

\begin{proof}
There are unique morphisms $u\colon I\to I'$ and $v\colon I'\to I$. Then $v\circ u$ and $\mathrm{id}_I$ are two morphisms $I\to I$, and there is only one, so $v\circ u=\mathrm{id}_I$; likewise $u\circ v=\mathrm{id}_{I'}$.
Uniqueness of the isomorphism is part of the hypothesis. Terminal objects of $\mathcal C$ are the initial objects of $\mathcal C^{\mathrm{op}}$.
\end{proof}

Almost every universal property says that some object is initial or terminal in a suitable auxiliary category, and Proposition~\ref{prop:initialunique} then gives its uniqueness for free.

\begin{definition}\label{def:product}
A \emph{product} of two objects $A,B$ is an object $P$ with morphisms $p_A\colon P\to A$, $p_B\colon P\to B$ such that for every object $X$ and morphisms $f\colon X\to A$, $g\colon X\to B$ there is exactly one
morphism $h\colon X\to P$ with $p_A\circ h=f$ and $p_B\circ h=g$ (Figure~\ref{fig:universal}). We write $P=A\times B$ and $h=(f,g)$. A \emph{coproduct} $A\sqcup B$ is a product in $\mathcal C^{\mathrm{op}}$: morphisms
$i_A\colon A\to A\sqcup B$, $i_B\colon B\to A\sqcup B$ such that every pair $f\colon A\to X$, $g\colon B\to X$ is $h\circ i_A$, $h\circ i_B$ for exactly one $h\colon A\sqcup B\to X$.
\end{definition}

A product of $A$ and $B$ is a terminal object in the category whose objects are triples $(X,f,g)$ and whose morphisms $(X,f,g)\to(X',f',g')$ are the $k\colon X\to X'$ with $f'\circ k=f$ and $g'\circ k=g$. So, by
Proposition~\ref{prop:initialunique}, \emph{the product is unique up to a unique isomorphism compatible with the projections}, and we may speak of \emph{the} product.

\begin{figure}[ht]
\centering
\fig{ch01-universal}
\caption{Left: the universal property of the product. Every pair of arrows $f$, $g$ out of $X$ factors through one arrow $h=(f,g)$ into $A\times B$. Right: the naturality square of a
natural transformation $\eta\colon F\Rightarrow G$. Going right then down equals going down then right.}
\label{fig:universal}
\end{figure}

\begin{example}\label{ex:products}
\leavevmode
\begin{enumerate}[label=(\alph*)]
\item In $\mathbf{Set}$, $\mathbf{Grp}$, $\mathbf{Ring}$ and $R\text{-}\mathbf{Mod}$ the product is the cartesian product with componentwise operations: $h(x)=(f(x),g(x))$ is the only possibility, and it is a
morphism.
\item In $\mathbf{Top}$ the product is $A\times B$ with the product topology, and the universal property explains the definition of that topology: it is the coarsest one making both projections
continuous, which is exactly what is needed for every $(f,g)$ with $f$, $g$ continuous to be continuous.
\item In a poset, $x\times y$ is the greatest lower bound $\inf(x,y)$, and $x\sqcup y$ the least upper bound, when they exist. In the poset of positive integers ordered by divisibility, they are the
$\gcd$ and the $\operatorname{lcm}$.
\item In $\mathbf{Set}$ the coproduct is the disjoint union. In $R\text{-}\mathbf{Mod}$ it is the direct sum $A\oplus B$ with $i_A(a)=(a,0)$, $i_B(b)=(0,b)$ and $h(a,b)=f(a)+g(b)$: the same object as the product,
a coincidence special to additive categories. In $\mathbf{Grp}$ the formula $h(a,b)=f(a)g(b)$ is not a homomorphism unless the images commute, and the coproduct is the free product $A*B$ of Tome~8,
Proposition~5.17, whose universal property is exactly the one required.
\end{enumerate}
\end{example}

A universal property can also characterise a single object with a map. Abelianisation is the first example.

\begin{proposition}[universal property of abelianisation]\label{prop:abuniv}
Let $G$ be a group. Every homomorphism $f\colon G\to A$ to an abelian group factors as $f=\bar f\circ\pi_G$ for exactly one homomorphism $\bar f\colon G^{\mathrm{ab}}\to A$.
\end{proposition}

\begin{proof}
Since $A$ is abelian,
\[f(aba^{-1}b^{-1})=f(a)f(b)f(a)^{-1}f(b)^{-1}=1,\]
so the kernel of $f$ contains every commutator, hence $G'$. By the factorisation theorem for quotients (Tome~2, Chapter~2), $\bar f(gG')=f(g)$ is a
well-defined homomorphism with $\bar f\circ\pi_G=f$. It is the only one, since $\pi_G$ is surjective and $\bar f$ is determined by its values on $\pi_G(G)=G^{\mathrm{ab}}$.
\end{proof}

In other words, $(G^{\mathrm{ab}},\pi_G)$ is initial among homomorphisms from $G$ to abelian groups: it is the best abelian approximation of $G$. Free objects are defined the same way: the free $R$-module on
a set $S$ is a module $R^{(S)}$ with a map $S\to R^{(S)}$ such that every map from $S$ to a module $M$ extends uniquely to a linear map $R^{(S)}\to M$; it is the module of finitely supported functions
$S\to R$, and the extension of $x\colon S\to M$ is $\sum_s c_s e_s\mapsto\sum_s c_s x(s)$.

\section{The Yoneda lemma}

Let $F\colon\mathcal C\to\mathbf{Set}$ be a functor and $A$ an object. Every element $x\in F(A)$ produces maps $\operatorname{Hom}(A,B)\to F(B)$, one for each $B$: send $f\colon A\to B$ to $F(f)(x)$. The Yoneda lemma
says that every natural transformation out of $h^A$ arises in this way, from exactly one $x$.

\begin{theorem}[Yoneda's lemma]\label{thm:yoneda}
Let $F\colon\mathcal C\to\mathbf{Set}$ be a functor and $A$ an object of $\mathcal C$. The map
\[\Phi\colon\operatorname{Nat}(h^A,F)\to F(A),\qquad\alpha\mapsto\alpha_A(\mathrm{id}_A),\]
is a bijection. Its inverse sends $x\in F(A)$ to the natural transformation $\alpha^x$ with $\alpha^x_B(f)=F(f)(x)$ for $f\colon A\to B$.
\end{theorem}

\begin{proof}
First, $\alpha^x$ is natural: for $g\colon B\to C$ and $f\in\operatorname{Hom}(A,B)$, $\alpha^x_C(h^A(g)(f))=F(g\circ f)(x)=F(g)(F(f)(x))=F(g)(\alpha^x_B(f))$, by functoriality of $F$. Next,
$\Phi(\alpha^x)=F(\mathrm{id}_A)(x)=x$. Finally, let $\alpha$ be any natural transformation and $x=\alpha_A(\mathrm{id}_A)$. For $f\colon A\to B$, naturality of $\alpha$ applied to $f$ and to the element $\mathrm{id}_A$
of $h^A(A)$ gives
\[\alpha_B(f)=\alpha_B(f\circ\mathrm{id}_A)=\alpha_B\bigl(h^A(f)(\mathrm{id}_A)\bigr)=F(f)\bigl(\alpha_A(\mathrm{id}_A)\bigr)=F(f)(x)=\alpha^x_B(f),\]
so $\alpha=\alpha^{\Phi(\alpha)}$. The two maps are inverse to each other.
\end{proof}

The whole proof rests on one observation: every morphism $f\colon A\to B$ is the image of $\mathrm{id}_A$ under $h^A(f)$, so a natural transformation out of $h^A$ is determined by what it does to $\mathrm{id}_A$.
Apply the lemma to $\mathcal C^{\mathrm{op}}$: for a contravariant $F$, natural transformations $h_A\Rightarrow F$ correspond to elements of $F(A)$, by the same formula.

\begin{corollary}[an object is known by its maps]\label{cor:yonedaemb}
For objects $A,B$, the natural transformations $h^A\Rightarrow h^B$ are exactly the maps $-\circ u$ for $u\in\operatorname{Hom}(B,A)$, and $h^A$ and $h^B$ are naturally isomorphic if and only if $A\cong B$. Likewise,
the natural transformations $h_A\Rightarrow h_B$ are the maps $u\circ-$ for $u\in\operatorname{Hom}(A,B)$, and $h_A\cong h_B$ if and only if $A\cong B$.
\end{corollary}

\begin{proof}
Take $F=h^B$ in Theorem~\ref{thm:yoneda}: $\operatorname{Nat}(h^A,h^B)$ is in bijection with $h^B(A)=\operatorname{Hom}(B,A)$, the transformation attached to $u$ being $f\mapsto h^B(f)(u)=f\circ u$. Under this bijection
composition of transformations corresponds to composition of morphisms in the opposite order: the transformation of $u$ followed by that of $v\colon C\to B$ sends $f$ to $f\circ u\circ v$, which is the transformation
of $u\circ v$. Inverse natural transformations therefore correspond to morphisms $u\colon B\to A$, $u'\colon A\to B$ with $u\circ u'=\mathrm{id}_A$ and $u'\circ u=\mathrm{id}_B$. The statement for $h_A$ is the same one in $\mathcal C^{\mathrm{op}}$.
\end{proof}

Two examples show that this abstract statement holds real content.

\begin{example}[Cayley's theorem]\label{ex:cayley}
Let $G$ be a group, seen as a category with one object $\bullet$. A functor $F\colon G\to\mathbf{Set}$ is a set $X=F(\bullet)$ with a map $F(g)\colon X\to X$ for each $g$, such that $F(gh)=F(g)F(h)$ and $F(1)=\mathrm{id}$: a $G$-set
(Tome~2, Chapter~3). A natural transformation between two $G$-sets is a map commuting with the actions. The functor $h^\bullet$ is $G$ acting on itself by left multiplication. By Corollary~\ref{cor:yonedaemb}, the maps
$G\to G$ commuting with all left multiplications are exactly the right multiplications $x\mapsto xu$, and Theorem~\ref{thm:yoneda} with $F=h^\bullet$ shows that distinct $g$ give distinct permutations
$x\mapsto gx$ of $G$. The second statement is Cayley's theorem: every group is a group of permutations of its own elements.
\end{example}

\begin{example}[the natural operations on groups]\label{ex:natops}
Which operations $G\to G$, $g\mapsto w_G(g)$, can be defined on every group, compatibly with all homomorphisms, that is, with $f(w_G(g))=w_H(f(g))$ for every $f\colon G\to H$? Such a family is exactly a natural
transformation $U\Rightarrow U$, where $U\colon\mathbf{Grp}\to\mathbf{Set}$ is the forgetful functor. Now $U$ is representable: $U\cong h^{\Z}$, since a homomorphism $f\colon\Z\to G$ is the same as the element $f(1)$. By Yoneda,
\[\operatorname{Nat}(U,U)\cong\operatorname{Nat}(h^{\Z},U)\cong U(\Z)=\Z.\]
The transformation attached to $n\in\Z$ sends $g=f(1)$ to $f(n)=g^n$. So the only operations on one element that make sense in every group are the powers $g\mapsto g^n$: there is no other formula, however
clever. Exercise~\ref{ex:words} does the same for operations on two elements.
\end{example}

\begin{think}
The forgetful functor $\mathbf{CRing}\to\mathbf{Set}$ is represented by the polynomial ring $\Z[X]$, because a ring homomorphism $\Z[X]\to R$ is determined by the image of $X$, which can be any element. What does Yoneda
say about the operations $R\to R$ that make sense in every commutative ring and commute with all ring homomorphisms?
\end{think}

\section{Adjunctions}

The free module, abelianisation and the cartesian product share one more pattern, which Daniel Kan isolated in 1958.

\begin{definition}\label{def:adjunction}
Functors $F\colon\mathcal C\to\mathcal D$ and $G\colon\mathcal D\to\mathcal C$ are \emph{adjoint}, $F$ \emph{left adjoint} to $G$ and $G$ \emph{right adjoint} to $F$, written $F\dashv G$, if there are bijections
\[\varphi_{A,B}\colon\operatorname{Hom}_{\mathcal D}(F(A),B)\to\operatorname{Hom}_{\mathcal C}(A,G(B)),\]
natural in $A$ and $B$: for $a\colon A'\to A$, $b\colon B\to B'$ and $h\colon F(A)\to B$, $\varphi_{A',B'}(b\circ h\circ F(a))=G(b)\circ\varphi_{A,B}(h)\circ a$.
\end{definition}

\begin{example}\label{ex:adjunctions}
\leavevmode
\begin{enumerate}[label=(\alph*)]
\item \emph{Free and forgetful.} The free module functor $S\mapsto R^{(S)}$ is left adjoint to the forgetful functor $U\colon R\text{-}\mathbf{Mod}\to\mathbf{Set}$: linear maps $R^{(S)}\to M$ correspond to maps $S\to U(M)$, by restriction to the
basis. The same holds for the free group on a set (Tome~8, Chapter~5) and the polynomial ring $\Z[X_s:s\in S]$, free commutative ring on $S$.
\item \emph{Abelianisation.} By Proposition~\ref{prop:abuniv}, $G\mapsto G^{\mathrm{ab}}$ is left adjoint to the inclusion $\mathbf{Ab}\to\mathbf{Grp}$: $\operatorname{Hom}_{\mathbf{Ab}}(G^{\mathrm{ab}},A)\cong\operatorname{Hom}_{\mathbf{Grp}}(G,A)$ by $\bar f\mapsto\bar f\circ\pi_G$.
\item \emph{Currying.} In $\mathbf{Set}$, for a fixed set $B$, maps $A\times B\to C$ correspond to maps $A\to C^B$ (the set of maps $B\to C$): $h\mapsto(a\mapsto h(a,-))$. So $-\times B\dashv(-)^B$. Programmers call this
\emph{currying}, after Haskell Curry.
\item \emph{Topologies.} The discrete topology $\mathbf{Set}\to\mathbf{Top}$ is left adjoint to the forgetful functor (every map out of a discrete space is continuous), and the indiscrete topology is right adjoint to it (every
map into an indiscrete space is continuous).
\item \emph{Galois connections.} Between posets, an adjunction $F\dashv G$ means $F(x)\le y\iff x\le G(y)$. For a finite Galois extension $L/K$ with group $\Gamma$, the maps $H\mapsto L^H$ from subgroups of $\Gamma$ to
intermediate fields, and $E\mapsto\operatorname{Gal}(L/E)$, both reverse inclusions, and $E\subseteq L^H\iff H\subseteq\operatorname{Gal}(L/E)$: an adjunction between the poset of subgroups and the opposite of the poset of
intermediate fields. The general notion is named after this example; the fundamental theorem of Galois theory (Tome~2, Chapter~10) says that here the adjunction is an equivalence.
\end{enumerate}
\end{example}

Adjoint functors are rigid: each determines the other, and a right adjoint inherits good behaviour from the bijection. The proofs are two applications of Yoneda.

\begin{theorem}[uniqueness of adjoints]\label{thm:adjunique}
If $G$ and $G'$ are both right adjoint to $F$, then $G\cong G'$: there is a natural isomorphism $G\Rightarrow G'$.
\end{theorem}

\begin{proof}
For each $B$, the composite bijections
\[\operatorname{Hom}(A,G(B))\cong\operatorname{Hom}(F(A),B)\cong\operatorname{Hom}(A,G'(B))\]
are natural in $A$, so they form a natural isomorphism $h_{G(B)}\Rightarrow h_{G'(B)}$. By
Corollary~\ref{cor:yonedaemb} it is $u_B\circ-$ for an isomorphism $u_B\colon G(B)\to G'(B)$, namely the image of $\mathrm{id}_{G(B)}$. For $b\colon B\to B'$, naturality of both bijections in $B$ shows that
$G'(b)\circ u_B$ and $u_{B'}\circ G(b)$ are both the image of $G(b)$ under the composite bijection for the pair $(G(B),B')$; hence they are equal, and $u$ is natural.
\end{proof}

\begin{theorem}[right adjoints preserve products]\label{thm:rapl}
Let $F\dashv G$ and let $B_1\times B_2$ be a product in $\mathcal D$, with projections $p_1,p_2$. Then $G(B_1\times B_2)$, with the projections $G(p_1),G(p_2)$, is a product of $G(B_1)$ and $G(B_2)$ in $\mathcal C$.
\end{theorem}

\begin{proof}
Let $A$ be an object of $\mathcal C$, with $f_1\colon A\to G(B_1)$ and $f_2\colon A\to G(B_2)$. Write $\psi=\varphi^{-1}$. By the universal property in $\mathcal D$ there is exactly one $k\colon F(A)\to B_1\times B_2$ with $p_j\circ k=\psi(f_j)$
for $j=1,2$. By naturality in $B$, $\varphi(p_j\circ k)=G(p_j)\circ\varphi(k)$, so $h=\varphi(k)$ satisfies $G(p_j)\circ h=f_j$. Conversely, if $h'\colon A\to G(B_1\times B_2)$ satisfies $G(p_j)\circ h'=f_j$, then
$k'=\psi(h')$ satisfies $\varphi(p_j\circ k')=G(p_j)\circ h'=f_j$, that is $p_j\circ k'=\psi(f_j)$, so $k'=k$ and $h'=h$.
\end{proof}

So the underlying set of a product of groups, rings, modules or spaces is the product of the underlying sets: all these forgetful functors have left adjoints. Dually, left adjoints preserve coproducts
(Exercise~\ref{ex:leftadj}), which explains at once why the free module on a disjoint union is a direct sum, and why the free group on $S\sqcup T$ is the free product of the free groups on $S$ and on $T$.
Abelianisation, a left adjoint, turns the coproduct of $\mathbf{Grp}$ into that of $\mathbf{Ab}$: $(A*B)^{\mathrm{ab}}\cong A^{\mathrm{ab}}\oplus B^{\mathrm{ab}}$, a fact one would otherwise prove by manipulating words.

\begin{beauty}
Mac Lane's slogan is that ``adjoint functors arise everywhere'', and this tome bears it out. In Chapter~2 the tensor product is left adjoint to $\operatorname{Hom}$, and this alone explains which exact sequences it preserves.
In Chapter~4, localisation is left adjoint to a forgetful functor. In Chapter~8, inducing a representation from a subgroup is right adjoint to restricting one (for finite groups it is left adjoint too), and this adjunction is Frobenius reciprocity, discovered by
Frobenius in 1898 as an identity between characters. Even computer science found the idea on its own: in 1991 Eugenio Moggi proposed to model the side effects of programs (exceptions, state, input and output) by
\emph{monads}, the structure that every adjunction $F\dashv G$ induces on the composite $G\circ F$, and Philip Wadler showed programmers how to use them; they became the way the language Haskell handles everything that is not a pure function.
The list type, with its ``flatten'', is the monad of the free monoid adjunction.
\end{beauty}

\section{Exercises}

\begin{exercise}[$\star$]\label{ex:zq}
Show that in $\mathbf{Ring}$ the inclusion $\Z\to\Q$ is a monomorphism and an epimorphism, but not an isomorphism. Is it an epimorphism in $\mathbf{Set}$?
\end{exercise}

\begin{exercise}[$\star$]\label{ex:nofieldproduct}
In the category of fields and ring homomorphisms, show that $\mathbb F_2$ and $\mathbb F_3$ have no product and no coproduct.
\end{exercise}

\begin{exercise}[$\star\star$]\label{ex:words}
Let $U\colon\mathbf{Grp}\to\mathbf{Set}$ be the forgetful functor and $U^2\colon G\mapsto G\times G$.
\begin{enumerate}[label=(\alph*)]
\item Show that $U^2$ is represented by the free group $F_2$ on two letters $a,b$, and deduce that the natural operations $G\times G\to G$ are exactly the maps $(g,h)\mapsto w(g,h)$ for a word $w\in F_2$.
\item Under the bijection $\operatorname{Nat}(U,U)\cong\Z$ of Example~\ref{ex:natops}, which integer is inversion? What does composition of natural operations become?
\end{enumerate}
\end{exercise}

\begin{exercise}[$\star$]\label{ex:matrixpoly}
Let $M_n\colon\mathbf{CRing}\to\mathbf{Set}$ send $R$ to the set of $n\times n$ matrices over $R$. Show that $M_n$ is represented by the polynomial ring $\Z[X_{ij}:1\le i,j\le n]$, and deduce that the families of maps
$M_n(R)\to R$ natural in $R$ are exactly the polynomials with integer coefficients in the entries. Name two of them.
\end{exercise}

\begin{exercise}[$\star\star$]\label{ex:leftadj}
\leavevmode
\begin{enumerate}[label=(\alph*)]
\item Prove that a left adjoint preserves coproducts. (Avoid work: use Theorem~\ref{thm:rapl} in the opposite categories.)
\item Deduce that $R^{(S\sqcup T)}\cong R^{(S)}\oplus R^{(T)}$, and that for sets $A\times(B\sqcup C)\cong(A\times B)\sqcup(A\times C)$.
\end{enumerate}
\end{exercise}

\begin{exercise}[$\star\star$]\label{ex:f2plane}
Let $K=\mathbb F_2$. For a plane $V$ (a two-dimensional space over $\mathbb F_2$) and $v,w\in V$, let $\eta_V(v)(w)=1$ if $v,w$ are linearly independent and $0$ otherwise.
\begin{enumerate}[label=(\alph*)]
\item Show that $\eta_V(v)$ is a linear form, that $\eta_V\colon V\to V^*$ is a linear isomorphism, and that \eqref{eq:dualnat} holds for every isomorphism $f$ between planes.
\item Where does the proof of Theorem~\ref{thm:notnatural} break down?
\end{enumerate}
\end{exercise}

\section{Hints}

\begin{hint}{ex:zq}
A ring homomorphism out of $\Q$ is determined by its values on $\Z$, because $\varphi(1/n)$ must be the inverse of $\varphi(n)$.
\end{hint}

\begin{hint}{ex:nofieldproduct}
A ring homomorphism between fields is injective and preserves the characteristic.
\end{hint}

\begin{hint}{ex:words}
A homomorphism $F_2\to G$ is determined by the images of $a$ and $b$, which can be any two elements.
\end{hint}

\begin{hint}{ex:matrixpoly}
A ring homomorphism $\Z[X_{ij}]\to R$ is the same as a choice of $n^2$ elements of $R$.
\end{hint}

\begin{hint}{ex:leftadj}
If $F\dashv G$ then $G^{\mathrm{op}}\dashv F^{\mathrm{op}}$ between the opposite categories, and a coproduct is a product in the opposite category.
\end{hint}

\begin{hint}{ex:f2plane}
A plane over $\mathbb F_2$ has three non-zero vectors, any two of which are independent.
\end{hint}

\section{Solutions}

\begin{solution}{ex:zq}
\idea Check the definitions with the hint: maps out of $\Q$ are already decided on $\Z$.
\writeup An injective map is a monomorphism in any concrete category (if $i\circ g=i\circ h$ then $g=h$ pointwise). Let $\varphi,\psi\colon\Q\to R$ agree on $\Z$. For $n\ne0$, $\varphi(n)\varphi(1/n)=\varphi(1)=1$, so $\varphi(n)$ is invertible and
$\varphi(m/n)=\varphi(m)\varphi(n)^{-1}$; the same formula holds for $\psi$, so $\varphi=\psi$. Hence the inclusion is an epimorphism. It is not an isomorphism, since there is no ring homomorphism $\Q\to\Z$ at all: $1/2$ would have to go to an inverse
of $2$ in $\Z$. In $\mathbf{Set}$ the epimorphisms are the surjections (if $f$ misses $y$, the maps $g,h$ to $\{0,1\}$ that differ only at $y$ satisfy $g\circ f=h\circ f$), so the inclusion is not an epimorphism there: being an epimorphism depends on the category.
\end{solution}

\begin{solution}{ex:nofieldproduct}
\idea A product needs maps to both fields, a coproduct maps from both, and a field has only one characteristic.
\writeup A ring homomorphism $\varphi\colon K\to L$ between fields is injective (its kernel is an ideal not containing $1$, hence $0$), so $K$ and $L$ have the same characteristic: $p\cdot1=0$ in $K$ if and only if it holds in $L$. A product
$P$ would come with homomorphisms $P\to\mathbb F_2$ and $P\to\mathbb F_3$, forcing $P$ to have characteristic $2$ and $3$. A coproduct $C$ would receive homomorphisms $\mathbb F_2\to C$ and $\mathbb F_3\to C$, with the same contradiction.
\end{solution}

\begin{solution}{ex:words}
\idea Repeat Example~\ref{ex:natops} with two generators.
\writeup (a) The map $\operatorname{Hom}(F_2,G)\to G\times G$, $f\mapsto(f(a),f(b))$, is a bijection by the universal property of the free group, and it is natural in $G$: for $u\colon G\to H$, $(u(f(a)),u(f(b)))$ is the image of $u\circ f$. So
$U^2\cong h^{F_2}$, and by Theorem~\ref{thm:yoneda}, $\operatorname{Nat}(U^2,U)\cong\operatorname{Nat}(h^{F_2},U)\cong U(F_2)=F_2$. The transformation attached to a word $w$ sends $(g,h)=(f(a),f(b))$ to $f(w)=w(g,h)$, the word with $g$ and $h$
substituted. For instance $ab$ gives the product, $aba^{-1}b^{-1}$ the commutator.
(b) Inversion is $g\mapsto g^{-1}$, the integer $-1$. The composite of $g\mapsto g^m$ and $g\mapsto g^n$ is $g\mapsto g^{mn}$: composition of natural operations is multiplication in $\Z$, so $\operatorname{Nat}(U,U)$ is the monoid $(\Z,\times)$.
\end{solution}

\begin{solution}{ex:matrixpoly}
\idea Same pattern: matrices are $n^2$ elements, and $n^2$ elements are a homomorphism out of the polynomial ring in $n^2$ variables.
\writeup Sending $f\colon\Z[X_{ij}]\to R$ to the matrix $(f(X_{ij}))$ is a bijection
\[
\operatorname{Hom}(\Z[X_{ij}],R)\to M_n(R),
\]
natural in $R$ (both sides transform by applying a homomorphism entrywise). So $M_n\cong h^{\Z[X_{ij}]}$ and
$\operatorname{Nat}(M_n,U)\cong U(\Z[X_{ij}])=\Z[X_{ij}]$, where $U$ is the forgetful functor. The transformation attached to a polynomial $P$ sends the matrix $(m_{ij})$ to $P(m_{ij})$. The determinant and the trace are two of them; so is the
coefficient of $t^k$ in the characteristic polynomial, for each $k$. This explains why identities like $\det(AB)=\det A\det B$ or Cayley--Hamilton, once proved over $\C$, hold over every commutative ring: both sides are natural, hence given by integer
polynomials, and polynomials that agree on all complex matrices are equal.
\end{solution}

\begin{solution}{ex:leftadj}
\idea Duality turns the statement into Theorem~\ref{thm:rapl}.
\writeup (a) Let $F\dashv G$ with bijections $\varphi_{A,B}$. Read in the opposite categories, the inverse bijections become
\[\operatorname{Hom}_{\mathcal C^{\mathrm{op}}}(G(B),A)\to\operatorname{Hom}_{\mathcal D^{\mathrm{op}}}(B,F(A)),\]
natural in both variables: $G^{\mathrm{op}}\dashv F^{\mathrm{op}}$. So $F^{\mathrm{op}}$ is a right adjoint and preserves products by Theorem~\ref{thm:rapl};
products in the opposite categories are coproducts. (b) The free module functor is a left adjoint (Example~\ref{ex:adjunctions}(a)), and the coproduct of sets is $S\sqcup T$, the coproduct of modules $\oplus$. In $\mathbf{Set}$, $A\times-$ is
left adjoint to $(-)^A$ by currying (with the factors swapped), so it maps the coproduct $B\sqcup C$ to the coproduct of $A\times B$ and $A\times C$.
\end{solution}

\begin{solution}{ex:f2plane}
\idea List everything: a plane over $\mathbb F_2$ has four vectors, $0,u,v,u+v$.
\writeup (a) Let $v\ne0$. The vectors independent of $v$ are exactly the two vectors outside $\{0,v\}$, so $\eta_V(v)$ is the function equal to $0$ on the line $\{0,v\}$ and $1$ off it: the unique non-zero linear form vanishing on $v$.
And $\eta_V(0)=0$. So $\eta_V$ sends the three non-zero vectors to the three non-zero forms, injectively (distinct lines have distinct kernels). It is additive: for distinct non-zero $v_1,v_2$, the form $\eta(v_1)+\eta(v_2)$ is non-zero (it takes the value
$1$ at $v_1$) and vanishes at $v_1+v_2$, since $\eta(v_1)(v_1+v_2)=\eta(v_1)(v_2)=1=\eta(v_2)(v_1)=\eta(v_2)(v_1+v_2)$; so it is $\eta(v_1+v_2)$. The remaining cases ($v_1=v_2$, or one of them zero) are immediate. Over $\mathbb F_2$
additivity is linearity, so $\eta_V$ is a linear isomorphism. Finally an isomorphism $f$ preserves linear independence, so $\eta_W(f(v))(f(w))=\eta_V(v)(w)$, which is \eqref{eq:dualnat}.
(b) The proof used a scalar $\lambda$ with $\lambda^2\ne1$; in $\mathbb F_2$ the only non-zero scalar is $1$. (In dimension $3$ and more, a longer argument shows that no natural isomorphism exists over any field.)
\end{solution}

\section*{Sources and further reading}
\addcontentsline{toc}{section}{Sources and further reading}

S.~Eilenberg and S.~Mac Lane, ``Group extensions and homology'', \emph{Ann.\ of Math.} 43 (1942), 757--831, and ``General theory of natural equivalences'', \emph{Trans.\ Amer.\ Math.\ Soc.} 58 (1945), 231--294. The story of the
Michigan lecture and Steenrod's solenoid is told by Mac Lane himself in H.~Bass, H.~Cartan, P.~Freyd, A.~Heller and S.~Mac Lane, ``Samuel Eilenberg 1913--1998'', \emph{Biographical Memoirs of the National Academy of
Sciences} 79 (2001). The meeting at the Gare du Nord is recalled by Y.~Kinoshita in his note on Yoneda's death, posted to the categories mailing list on 23 April 1996. D.~M.~Kan, ``Adjoint functors'', \emph{Trans.\ Amer.\ Math.\ Soc.} 87 (1958),
294--329.

E.~Moggi, ``Notions of computation and monads'', \emph{Information and Computation} 93 (1991), 55--92; P.~Wadler, ``Monads for functional programming'', in \emph{Advanced Functional Programming}, Lecture Notes in Computer
Science 925 (Springer, 1995).

For more: S.~Mac Lane, \emph{Categories for the Working Mathematician} (Springer, 2nd ed.\ 1998), Chapters~I--IV; T.~Leinster, \emph{Basic Category Theory} (Cambridge University Press, 2014), short and very clear, freely available
from its author; E.~Riehl, \emph{Category Theory in Context} (Dover, 2016), with examples from every part of mathematics.
